%% 8. DISCUSSION
%% ============================================================================
%% 8.0 Benchmark retraction (2026-09-19)
%% ---------------------------------------------------------------------
\label{sec:discussion-retraction}
\paragraph{Benchmark retraction (2026-09-19).}
The 8-cell headline aggregates (Table~\ref{tab:headline_results})
were retracted on \texttt{audit-hostile/2026-09-19-fixes} after a
hostile review surfaced three data-integrity issues in the
iteration-1 draft runs: (i) the T5 and T7 rows were byte-identical
across all recorded metrics, which is physically impossible for two
distinct end-to-end LLM agent runs; (ii) the P0 winner was recorded
against the pre-registered decision rule, which (gates
$\rightarrow$ impl\_q $\rightarrow$ plan\_q $\rightarrow$ cost)
selects Setup A on impl quality 3.5 vs 2.0, ~4$\times$ fewer tokens,
and ~2.5$\times$ lower cost; (iii) two latency tables on the same
page contradicted each other by 484 seconds per run. The surviving
5-cell evidence (Table~\ref{tab:per_task_verdicts_surviving},
Table~\ref{tab:headline_results_v2}) yields an exact two-sided
sign-test p-value of 1.00 (coin flip); we therefore make no
directional product claim and do not present the n=1 per cell data
as evidence. The pre-registered iter-2 acceptance criterion
($\geq$7/8 wins with paired bootstrap 95\% CI excluding 0 on impl
quality) requires additional n$\geq$3 evidence before any directional
claim can be made. The retraction commit is preserved in the audit
trail so the original draft numbers remain reproducible.


%% 8.1 Key Findings

The historical controlled study reported the following directional findings
for one codebase and an earlier risk-label schema. They should not be read as
current real-usage or causal product evidence:

\begin{itemize}
    \item PLAN QUALITY SIGNAL: The sidecar condition received higher historical plan-quality and repo-fit scores (+0.32 and +0.43), subject to the study's task set and judge calibration.
    \item DIFF DISCIPLINE SIGNAL: The sidecar condition received a higher historical diff-discipline score (+0.23); this is not yet a validated reduction in harmful real-session edits.
    \item TOKEN EFFICIENCY SIGNAL: The study measured 8.2 percent lower total tokens on average, with up to 29 percent savings on cache-heavy tasks; current real-usage savings are unproven.
    \item CORRECTNESS SIGNAL: The study observed 100 percent pass rates on locked and rotated tests for the sidecar arm (vs 92 percent/90 percent for vanilla); this has not been reproduced under the current schema and default posture.
    \item LOCAL OPERATION: The entire system runs locally with no external telemetry, addressing privacy concerns.
\end{itemize}

%% 8.2 Performance Trade-offs

The system introduces a performance trade-off:

\begin{itemize}
    \item ADDITIONAL LATENCY: The sidecar adds approximately 98 seconds per run due to structural analysis overhead.
    \item MEMORY USAGE: The Mamba-130M model requires ~200MB RAM.
    \item CPU USAGE: The system uses CPU for inference.
\end{itemize}

These costs may have been offset in the historical task set, but that trade-off
has not been established in current repository sessions:
\begin{itemize}
    \item HISTORICAL TOKEN SIGNAL: 8.2 percent average token savings in the controlled study.
    \item UNVALIDATED REGRESSION HYPOTHESIS: Fewer harmful or off-plan edits may reduce manual intervention.
    \item UNVALIDATED QUALITY HYPOTHESIS: Better code may require less review and revision.
\end{itemize}

%% 8.3 Why This Works

The historical result patterns are consistent with several design decisions,
but the contribution of each decision has not been isolated:

\begin{itemize}
    \item HYBRID ARCHITECTURE: Combining System 1 (LLM) and System 2 (structural specialist)
    \item LIGHTWEIGHT SPECIALIST: Mamba-130M is small enough to run locally
    \item TREE-SITTER PARSING: Using a proven, fast parser for AST extraction
    \item PER-FILE-KIND THRESHOLDS: Different file types have different acceptable risk profiles
    \item COLD-START HANDLING: New files are handled gracefully
    \item SHADOW MODE: Allows safe observation before enforcement
\end{itemize}

%% 8.4 Comparison to Alternative Approaches

\begin{description}
    \item[PURE LLM SCALING] Increasing model size and context window does not address the structural blind spot.
    \item[TOOL-USING AGENTS] Agents that call external tools can catch some structural issues, but these calls are stateless.
    \item[RETRIEVAL-AUGMENTED GENERATION] RAG can provide relevant context, but does not inherently understand code structure.
    \item[FINE-TUNED MODELS] Domain-specific fine-tuning can improve code understanding, but requires significant resources.
\end{description}

Our hybrid approach combines the strengths of these approaches while addressing their limitations.

%% PLACEHOLDER: Add comparison to specific alternative systems
%% PLACEHOLDER: Add discussion of theoretical implications
%% PLACEHOLDER: Add implications for AI safety and alignment
